\subsubsection{Recourse-Aware Mechanism}
\label{sec:recourse_aware_mechanism}

Each decision epoch contains two sequential decision states. The pre-recourse state $S_t^1$ is observed before the human-driven electric vehicle (HEV) assignment, whereas the realized residual state $\bar S_t^2$ is observed only after the HEV acceptance and rejection outcomes are known. The phase-augmented state space and within-epoch transition are
\begin{equation}\mathcal S=\mathcal S^1\mathbin{\dot\cup}\mathcal S^2,\qquad S_t^1\xrightarrow{\,X_t\,}\bar S_t^2\xrightarrow{\,Y_t\,}S_{t+1}^1.\label{eq:recourse_state_sequence}\end{equation}

Let $h\in\mathcal V_t^{\mathrm{HEV}}$ and $m\in\mathcal V_t^{\mathrm{AEV}}$ index an HEV and an autonomous electric vehicle (AEV), respectively. In the first stage, $x_{jht}=1$ means that request $j$ is offered to HEV $h$, while $n_{ht}=1$ sends an HEV receiving no request to the driver-side relocation process. Collecting these variables in $\boldsymbol x_t$ and $\boldsymbol n_t$ gives
\begin{equation}X_t=(\boldsymbol x_t,\boldsymbol n_t)\in\mathcal F_t^1(S_t^1).\label{eq:recourse_stage1_action}\end{equation}

After $X_t$ is executed, the random outcome $\xi_t$ partitions the affected requests into accepted, rejected, and unassigned sets, denoted by $\mathcal J_t^{\mathrm{acc}}$, $\mathcal J_t^{\mathrm{rej}}$, and $\mathcal J_t^{\mathrm{rem}}$. Accepted requests leave the available-request pool, and the realized residual state is
\begin{equation}\bar S_t^2=\Phi_t(S_t^1,X_t,\xi_t).\label{eq:recourse_residual_state}\end{equation}
The AEV stage operates on the residual request set; a rejected request becomes a recourse request only if it is actually assigned to an AEV:
\begin{equation}\mathcal J_t^2=\mathcal J_t^{\mathrm{rej}}\mathbin{\dot\cup}\mathcal J_t^{\mathrm{rem}},\qquad \mathcal J_t^{\mathrm{rec}}(Y_t)=\left\{j\in\mathcal J_t^{\mathrm{rej}}:\sum_{m\in\mathcal V_t^{\mathrm{AEV}}}u_{jmt}=1\right\}.\label{eq:recourse_request_sets}\end{equation}

The realized first-stage reward contains accepted-service, rejection, and HEV-relocation components:
\begin{equation}R_t^1=R_t^{1,\mathrm{acc}}+R_t^{1,\mathrm{rej}}+R_t^{1,\mathrm{rel}}.\label{eq:recourse_stage1_reward}\end{equation}

In the second stage, $u_{jmt}$, $c_{mct}$, $a_{mzt}$, and $w_{mt}$ indicate that AEV $m$ serves request $j$, charges at station $c$, relocates to zone $z$, or waits, respectively. Collecting the four decision arrays in $\boldsymbol u_t$, $\boldsymbol c_t$, $\boldsymbol a_t$, and $\boldsymbol w_t$ gives
\begin{equation}Y_t=(\boldsymbol u_t,\boldsymbol c_t,\boldsymbol a_t,\boldsymbol w_t)\in\mathcal F_t^2(\bar S_t^2).\label{eq:recourse_stage2_action}\end{equation}
Its reward contains recourse service, service of previously unassigned requests, charging, relocation, and waiting; the recourse term is part of the AEV reward and is not added again elsewhere:
\begin{equation}R_t^2=R_t^{2,\mathrm{rec}}+R_t^{2,\mathrm{rem}}+R_t^{2,\mathrm{chg}}+R_t^{2,\mathrm{rel}}+R_t^{2,\mathrm{wait}}.\label{eq:recourse_stage2_reward}\end{equation}

For policy $\pi=(\pi_1,\pi_2)$, the first-stage action value backs up the value of the residual state generated by $X_t$:
\begin{equation}Q_{1,t}^{\pi}(S_t^1,X_t)=\mathbb E^{\pi}\!\left[R_t^1+Q_{2,t}^{\pi}\!\left(\bar S_t^2,\pi_2(\bar S_t^2)\right)\mid S_t^1,X_t\right],\qquad \bar S_t^2=\Phi_t(S_t^1,X_t,\xi_t).\label{eq:recourse_q1}\end{equation}
The second-stage action value contains the transition to the next decision epoch:
\begin{equation}Q_{2,t}^{\pi}(\bar S_t^2,Y_t)=\mathbb E^{\pi}\!\left[R_t^2+\gamma^{\Delta t_t}Q_{1,t+1}^{\pi}\!\left(S_{t+1}^1,\pi_1(S_{t+1}^1)\right)\mid \bar S_t^2,Y_t\right].\label{eq:recourse_q2}\end{equation}
No discount is applied between the two stages of the same epoch. Substitution gives the closed stage-1 recursion
\begin{equation}Q_{1,t}^{\pi}(S_t^1,X_t)=\mathbb E^{\pi}\!\left[R_t^1+R_t^2+\gamma^{\Delta t_t}Q_{1,t+1}^{\pi}\!\left(S_{t+1}^1,\pi_1(S_{t+1}^1)\right)\mid S_t^1,X_t\right].\label{eq:recourse_q1_expanded}\end{equation}

The central distinction from independently evaluating two stage-specific states is the endogenous residual state $\bar S_t^2$: it is generated by the first-stage decision and its realized rejection outcome, rather than supplied as an unrelated second state. Consequently, $Q_{1,t}^{\pi}$ learns the expected AEV recourse value through $Q_{2,t}^{\pi}$, including the value of serving HEV-rejected requests. The two value functions therefore form a conditional recursion, not two parallel values to be computed independently or added after optimization. Feasibility and cross-vehicle coupling remain enforced by $\mathcal F_t^1(S_t^1)$ and $\mathcal F_t^2(\bar S_t^2)$.

The corresponding stage policies are
\begin{equation}\pi_1(S_t^1)\in\operatorname*{arg\,max}_{X_t\in\mathcal F_t^1(S_t^1)}Q_{1,t}^{\pi}(S_t^1,X_t).\label{eq:recourse_policy1}\end{equation}
\begin{equation}\pi_2(\bar S_t^2)\in\operatorname*{arg\,max}_{Y_t\in\mathcal F_t^2(\bar S_t^2)}Q_{2,t}^{\pi}(\bar S_t^2,Y_t).\label{eq:recourse_policy2}\end{equation}
